\section{From Personal Inventory to Opportunity Discovery}

\subsection{Where Assets Collide with Markets}

A list of skills is not yet a business. A list of personal frustrations is not yet a business either. The interesting work begins when the 2 meet: what problems can you see, and what can you do about them that a buyer may value?

Paul Graham's essay \textit{How to Get Startup Ideas} argues for noticing problems, particularly ones you understand through direct experience, rather than forcing yourself to invent concepts in the abstract.\footnote{Paul Graham, ``How to Get Startup Ideas'' (November 2012), \url{https://www.paulgraham.com/startupideas.html}.} This is a useful antidote to trend chasing. But it needs an important qualification: a problem you experience may be idiosyncratic, trivial to other buyers, or something others are unwilling to pay to solve. \textbf{Founder insight improves the search space; it does not complete market validation.}

There are 6 productive routes to generate opportunities from starting assets:

\subsection{6 Discovery Routes}

\textbf{Route A: Work people already ask you to do.} If multiple people seek your help for the same class of problem, you possess a clue about a potentially valuable service. Ask: Who repeatedly needs this? Can the result be delivered and priced? Can someone besides you deliver it? \\
\textit{Danger:} Free requests from friends often reflect zero price rather than commercial willingness to pay.

\textbf{Route B: Costly frustrations you have witnessed.} Look for repeated delays, errors, rework, avoidable risk, or broken workflows. An engineer may see machines fail because maintenance logs are never analyzed; an accountant may see reconciliation tasks consume days each month. \\
\textit{Danger:} The person suffering the pain may not control the budget. You must identify the economic buyer.

\textbf{Route C: Underused capabilities.} When a starting capability is exceptionally strong, search for buyers who value it. A research statistician may have experimental design skills scarce among marketing agencies. \\
\textit{Danger:} Searching for applications of a favorite hammer easily turns into a solution looking for a problem.

\textbf{Route D: Trusted customer access.} You may know school principals, plant managers, or developer community leads. The goal is not to pitch them blindly, but to learn what problems they already allocate capital and attention to solve. \\
\textit{Danger:} Access without delivery capability tempts founders to overpromise before securing execution.

\textbf{Route E: Environmental shifts.} When technology, regulations, or input costs change, longstanding problems suddenly become solvable at a fraction of their prior cost. \\
\textit{Danger:} A powerful new technology does not guarantee a willing buyer or defensible unit economics.

\textbf{Route F: Unique intersections.} Some founders possess combinations that are uncommon even if neither skill is world-class: software plus clinical operations, statistics plus teaching, or logistics plus local procurement. \textbf{Intersections are fertile because they let you translate between groups that misunderstand each other, producing outcomes existing silos cannot coordinate to deliver.} \\
\textit{Danger:} An unusual combination is not automatically valuable; customers must have urgent need for the bridge.
